% 康奈尔听课笔记模板（配色取自王道课件）
\documentclass[10.5pt,a4paper,oneside]{article}
\usepackage[a4paper,top=2.3cm,bottom=2.2cm,left=5.7cm,right=1.5cm,
  marginparwidth=4.25cm,marginparsep=0.55cm,headheight=14pt,footskip=1.1cm]{geometry}
\usepackage{xeCJK}
\usepackage{xeCJKfntef}
\xeCJKDeclareCharClass{FullLeft}{"201C,"2018}
\xeCJKDeclareCharClass{FullRight}{"201D,"2019}
\xeCJKDeclareCharClass{CJK}{"2460->"24FF,"2605,"2606}
\setmainfont{Noto Sans CJK SC}[BoldFont={Noto Sans CJK SC Bold},ItalicFont={Noto Sans CJK SC},ItalicFeatures={FakeSlant=0.18}]
\setsansfont{Noto Sans CJK SC}[BoldFont={Noto Sans CJK SC Bold}]
\setmonofont{DejaVu Sans Mono}[Scale=0.9]
\setCJKmainfont{Noto Sans CJK SC}[BoldFont={Noto Sans CJK SC Bold},ItalicFont={Noto Sans CJK SC},ItalicFeatures={FakeSlant=0.18}]
\setCJKsansfont{Noto Sans CJK SC}[BoldFont={Noto Sans CJK SC Bold}]
\setCJKmonofont{Noto Sans Mono CJK SC}[Scale=0.9]
\linespread{1.32}
\usepackage{amsmath,amssymb}
\usepackage[dvipsnames,table]{xcolor}
\usepackage{graphicx}
\usepackage{adjustbox}
\usepackage{tikz}
\usetikzlibrary{calc}
\usepackage[most]{tcolorbox}
\usepackage{eso-pic}
\usepackage{fancyhdr}
\usepackage{titlesec}
\usepackage{needspace}
\usepackage{enumitem}
\usepackage{xltabular,booktabs,array,colortbl}
\usepackage{marginnote}
\reversemarginpar
\renewcommand*{\raggedleftmarginnote}{}
\renewcommand*{\raggedrightmarginnote}{}
\usepackage{caption}
\usepackage{titletoc}
\usepackage{fancyvrb}
\usepackage{framed}
\usepackage[section]{placeins}
\usepackage{hyperref}
\hypersetup{colorlinks=true,linkcolor=wdblue,urlcolor=wdblue,pdftitle={3.4.1-3.4.4
特殊矩阵的压缩存储},pdfauthor={听课笔记}}

% ---------- 课件配色 ----------
\definecolor{wdblue}{HTML}{4472C4}
\definecolor{wdlightblue}{HTML}{5B9BD5}
\definecolor{wdtitle}{HTML}{4A86C8}
\definecolor{wddark}{HTML}{1F3864}
\definecolor{wdorange}{HTML}{ED7D31}
\definecolor{wdgray}{HTML}{A5A5A5}
\definecolor{wdgold}{HTML}{FFC000}
\definecolor{wdgreen}{HTML}{70AD47}
\definecolor{wddeepgreen}{HTML}{3F7F2A}
\definecolor{wdred}{HTML}{C00000}
\definecolor{wdpurple}{HTML}{7030A0}
\definecolor{cuebg}{HTML}{F2F6FC}
\definecolor{codebg}{HTML}{F7F9FC}
\color[HTML]{262626}

% ---------- 行内样式 ----------
\newcommand{\hlred}[1]{\textcolor{wdred}{\bfseries #1}}
\newcommand{\hlblue}[1]{\textcolor{wdblue}{\bfseries #1}}
\newcommand{\hlgreen}[1]{\textcolor{wddeepgreen}{\bfseries #1}}
\newcommand{\hlorange}[1]{\textcolor{wdorange}{\bfseries #1}}
\newcommand{\hlpurple}[1]{\textcolor{wdpurple}{\bfseries #1}}
\newcommand{\uln}[1]{\CJKunderline[format=\color{wdred},thickness=0.9pt,depth=0.28em]{#1}}
\newcommand{\wavy}[1]{\CJKunderwave[format=\color{wdorange}]{#1}}
\newcommand{\hlmark}[1]{\CJKunderline[format=\color{wdgold!55},thickness=0.55em,depth=-0.05em]{#1}}
\newcommand{\tagbox}[2]{\tcbox[on line,colback=#1,colframe=#1,coltext=white,boxsep=0.5pt,left=3pt,right=3pt,top=0.5pt,bottom=0.5pt,arc=2pt,fontupper=\small\bfseries]{#2}}

% ---------- 康奈尔版式背景：左侧线索栏 ----------
\newlength{\cuestripw}\setlength{\cuestripw}{5.15cm}
\AddToShipoutPictureBG{%
  \begin{tikzpicture}[remember picture,overlay]
    \fill[cuebg] (current page.north west) rectangle ($(current page.south west)+(\cuestripw,0)$);
    \draw[wdblue,line width=1.2pt] ($(current page.north west)+(\cuestripw,-1.55cm)$) -- ($(current page.south west)+(\cuestripw,1.55cm)$);
    \node[anchor=west,text=wdblue,font=\small\bfseries] at ($(current page.north west)+(0.75cm,-1.25cm)$) {线索 · 问题};
    \node[anchor=west,text=wdblue,font=\small\bfseries] at ($(current page.north west)+(\cuestripw+0.55cm,-1.25cm)$) {笔记 · 老师主线 + 补充讲解};
  \end{tikzpicture}}

% ---------- 页眉页脚 ----------
\pagestyle{fancy}\fancyhf{}
\renewcommand{\headrulewidth}{0pt}
\fancyhead[R]{\small\color{wdgray}第3章 栈、队列和数组}
\fancyfoot[C]{\small\color{wdgray}\thepage}
\fancyfoot[R]{\small\color{wdorange}3.4.1-3.4.4
特殊矩阵的压缩存储 · 听课笔记}
\fancyfoot[L]{\small\hyperlink{toc}{\textcolor{wdblue}{\ensuremath{\uparrow}\,回到目录}}}
\fancypagestyle{plain}{\fancyhf{}\fancyfoot[C]{\small\color{wdgray}\thepage}\fancyfoot[R]{\small\color{wdorange}3.4.1-3.4.4
特殊矩阵的压缩存储 · 听课笔记}}

% ---------- 标题：仿课件的蓝色标题 + 六色条 ----------
\newcommand{\sixbar}{\par\nointerlineskip\vspace{3pt}\centering
  \begin{tikzpicture}\foreach \c [count=\i] in {wdlightblue,wdorange,wdgray,wdgold,wdblue,wdgreen}
     {\fill[\c] ({(\i-1)*1.05},0) rectangle ({(\i-1)*1.05+0.92},0.11);}\end{tikzpicture}\par}
\titleformat{\section}[block]{\centering\Large\bfseries\color{wdtitle}}{}{0pt}{}[\sixbar]
\titlespacing*{\section}{0pt}{16pt}{10pt}
\pretocmd{\section}{\Needspace{8\baselineskip}}{}{}
\titleformat{\subsection}{\Needspace{7\baselineskip}\large\bfseries\color{wddark}}{\textcolor{wdblue}{\rule[-0.15em]{4pt}{1.05em}}\hspace{6pt}}{0pt}{}
\titlespacing*{\subsection}{0pt}{12pt}{5pt}
\titleformat{\subsubsection}{\Needspace{5\baselineskip}\normalsize\bfseries\color{wdblue}}{\textcolor{wdorange}{\large\ensuremath{\blacktriangleright}}\hspace{4pt}}{0pt}{}
\titlespacing*{\subsubsection}{0pt}{8pt}{3pt}
\setcounter{secnumdepth}{0}
\setlength{\parindent}{0pt}\setlength{\parskip}{5pt}
\setlist{nosep,leftmargin=1.5em,itemsep=2pt}
\setlist[itemize,1]{label=\textcolor{wdblue}{\small\ensuremath{\bullet}}}
\setlist[itemize,2]{label=\textcolor{wdorange}{\scriptsize\ensuremath{\blacktriangleright}}}
\setlist[enumerate,1]{label=\textcolor{wdblue}{\bfseries\arabic*.}}
\providecommand{\tightlist}{\setlength{\itemsep}{2pt}\setlength{\parskip}{0pt}}

% ---------- 目录 ----------
\setcounter{tocdepth}{1}
\titlecontents{section}[0pt]{\addvspace{5pt}\bfseries\color{wddark}}{}{}{\titlerule*[0.7pc]{\color{wdgray!60}.}\contentspage}
\titlecontents{subsection}[1.6em]{\addvspace{1pt}\small\color{black!75}}{}{}{\titlerule*[0.7pc]{\color{wdgray!40}.}\contentspage}
\newcommand{\tocback}{\par\vspace{-3pt}{\hfill\footnotesize\hyperlink{toc}{\textcolor{wdgray}{\ensuremath{\uparrow}\,回到目录}}}\par\vspace{-4pt}}

% ---------- 线索栏 ----------
\newcommand{\cueblock}[2]{\Needspace{4\baselineskip}\leavevmode#1{#2}\par\vspace{-\baselineskip}\vspace{-\parskip}}
% 线索直接贴在段首行的左侧（不用 marginnote，避免分页时偶尔落进正文栏）
\newcommand{\cuebox}[1]{\leavevmode\llap{\smash{\begin{minipage}[t]{\marginparwidth}\raggedright\small\linespread{1.18}\selectfont #1\end{minipage}}\hspace{\marginparsep}}}
\newcommand{\cuenote}[1]{\cuebox{{\color{wdblue}\bfseries Q\,}\color{wddark}#1}}
\newcommand{\cuekey}[1]{\cuebox{\color{wdred}\bfseries #1}}

% ---------- 讲解框 ----------
\tcbset{wdbox/.style={enhanced,breakable,arc=2pt,boxrule=0pt,leftrule=3pt,left=7pt,right=7pt,top=5pt,bottom=5pt,
   before skip=7pt,after skip=7pt,fonttitle=\bfseries\small,attach title to upper={\quad},coltitle=#1,
   colframe=#1,colback=#1!6}}
\newtcolorbox{defbox}[1][]{enhanced,breakable,arc=3pt,boxrule=0.8pt,colframe=wdblue,colback=wdblue!4,
   left=8pt,right=8pt,top=6pt,bottom=6pt,fonttitle=\bfseries\small,coltitle=white,colbacktitle=wdblue,
   title={定义#1},before skip=8pt,after skip=8pt}
\newtcolorbox{thmbox}[1]{enhanced,breakable,arc=3pt,boxrule=0.8pt,colframe=wdblue,colback=wdblue!4,
   left=8pt,right=8pt,top=6pt,bottom=6pt,fonttitle=\bfseries\small,coltitle=white,colbacktitle=wdblue,
   title={#1},before skip=8pt,after skip=8pt}
\newtcolorbox{bridgebox}{wdbox=wdblue,title={承上启下}}
\newtcolorbox{plainbox}{wdbox=wddeepgreen,colframe=wdgreen,colback=wdgreen!8,title={大白话}}
\newtcolorbox{deepbox}{wdbox=wdpurple,title={背后的思想}}
\newtcolorbox{exambox}{wdbox=wdred,title={考点}}
\newtcolorbox{pitfallbox}{wdbox=wdorange,title={易错提醒}}
\newtcolorbox{bookbox}{wdbox=wddark,colback=wdgray!10,colframe=wdgray,title={王道书补充}}
\newtcolorbox{teacherbox}{wdbox=wdorange,colback=wdgold!8,colframe=wdgold,coltitle=wdorange,title={老师原话}}
\newtcolorbox{fixbox}{wdbox=wdred,colback=wdred!4,title={勘误 / 说明}}
% 习题笔记用
\newtcolorbox{questionbox}[1][]{enhanced,breakable,arc=3pt,boxrule=0.8pt,colframe=wddark,colback=white,
   left=8pt,right=8pt,top=6pt,bottom=6pt,fonttitle=\bfseries\small,coltitle=white,colbacktitle=wddark,
   title={#1},before skip=10pt,after skip=6pt}
\newtcolorbox{videobox}{wdbox=wdorange,colback=wdgold!8,colframe=wdgold,coltitle=wdorange,title={视频讲解}}
\newtcolorbox{jiexibox}{wdbox=wddark,colback=wdgray!10,colframe=wdgray,title={书上解析}}
\newtcolorbox{addbox}{wdbox=wdpurple,title={我的补充}}
\newtcolorbox{recallbox}{wdbox=wdblue,title={回顾正课}}
\newtcolorbox{watchbox}{enhanced,breakable,arc=3pt,boxrule=1pt,colframe=wdorange,colback=wdgold!12,
   fonttitle=\bfseries\small,coltitle=white,colbacktitle=wdorange,title={建议去看视频},left=8pt,right=8pt,before skip=8pt,after skip=8pt}
\newtcolorbox{summarybox}{enhanced,arc=3pt,
   boxrule=1pt,colframe=wdblue,colback=white,colbacktitle=wdblue,coltitle=white,fonttitle=\bfseries,
   title={总结区 · 用自己的话复述这一节},left=10pt,right=10pt,top=8pt,bottom=8pt,before skip=14pt}

% ---------- 代码 ----------

\makeatletter
\@ifundefined{Shaded}{}{\renewenvironment{Shaded}{\def\FrameCommand{\fboxsep=5pt\fboxrule=0.5pt\fcolorbox{wdblue!35}{codebg}}%
   \MakeFramed{\advance\hsize-\width\FrameRestore}\small}{\endMakeFramed}}
\makeatother
\fvset{fontsize=\small}

% ---------- 表格 ----------
\renewcommand{\arraystretch}{1.42}
\newcolumntype{Y}{>{\raggedright\arraybackslash}X}

% ---------- 图片 ----------
% 课件截图做成浮动体：LaTeX 会自动用正文填满页面，避免大块空白
\setcounter{topnumber}{3}\setcounter{bottomnumber}{2}\setcounter{totalnumber}{4}
\renewcommand{\topfraction}{0.85}\renewcommand{\bottomfraction}{0.6}\renewcommand{\textfraction}{0.12}\renewcommand{\floatpagefraction}{0.75}
\setlength{\textfloatsep}{10pt plus 3pt minus 2pt}\setlength{\floatsep}{8pt plus 2pt}
\newcommand{\slidefig}[3]{\begin{figure}[!ht]\centering\begin{tcolorbox}[enhanced,width=#2\linewidth,arc=2pt,boxrule=0.6pt,colframe=wdgray!70,
   colback=white,left=2pt,right=2pt,top=2pt,bottom=2pt,drop fuzzy shadow=black!25]\includegraphics[width=\linewidth]{#1}\end{tcolorbox}
   \vspace{-3pt}{\small\color{wdgray}#3}\end{figure}}
\newcommand{\plainfig}[3]{\par\noindent\begin{minipage}{\linewidth}\centering\includegraphics[width=#2\linewidth]{#1}\par\vspace{2pt}{\small\color{wdgray}#3}\end{minipage}\par\vspace{4pt}}

\begin{document}
% ---------- 封面条（仿课件封面） ----------
\thispagestyle{plain}
\begin{tcolorbox}[enhanced,enlarge left by=-4.8cm,width=\linewidth+4.8cm,arc=0pt,boxrule=0pt,
   colback=wddark,interior style={left color=wddark,right color=wdblue},left=18pt,right=18pt,top=14pt,bottom=12pt,
   overlay={\draw[white,line width=2.5pt] ([xshift=10pt,yshift=-10pt]frame.north west) -- ++(0,-0.9cm);
            \draw[white,line width=2.5pt] ([xshift=10pt,yshift=-10pt]frame.north west) -- ++(1.6cm,0);
            \draw[white,line width=2.5pt] ([xshift=-10pt,yshift=10pt]frame.south east) -- ++(0,0.9cm);
            \draw[white,line width=2.5pt] ([xshift=-10pt,yshift=10pt]frame.south east) -- ++(-1.6cm,0);}]
  {\color{white!85!wdlightblue}\small 数据结构 · 第3章
栈、队列和数组 · 康奈尔听课笔记}\par\vspace{6pt}
  {\color{white}\fontsize{22}{28}\selectfont\bfseries 3.4.1-3.4.4
特殊矩阵的压缩存储}\par\vspace{6pt}
  {\color{white!80!wdlightblue}\small 课件 3.4.1～3.4.4 · 正课字幕 P32 ·
王道书 §3.4.1--3.4.4（第 100--104 页）}
\end{tcolorbox}
\vspace{4pt}
\begin{tcolorbox}[enhanced,arc=3pt,boxrule=0.8pt,colframe=wdblue,colback=wdblue!3,left=10pt,right=10pt,top=6pt,bottom=6pt,
   colbacktitle=wdblue,coltitle=white,fonttitle=\bfseries\small,title={目录 · 点击标题跳转}]
\hypertarget{toc}{}\makeatletter\@starttoc{toc}\makeatother
\end{tcolorbox}
\hypertarget{ux8fd9ux4e00ux8282ux5728ux8bb2ux4ec0ux4e48}{%
\section{这一节在讲什么}\label{ux8fd9ux4e00ux8282ux5728ux8bb2ux4ec0ux4e48}}

\tocback

\Needspace{4\baselineskip}

\cuenote{这一节讲哪几种矩阵？所有压缩存储题的共同做法是什么？}矩阵可以很方便地用数组存储，所以老师先讲一维、二维数组背后的实现细节，再讲四种特殊矩阵：\textbf{对称矩阵}、\textbf{三角矩阵}、\textbf{三对角矩阵}、\textbf{稀疏矩阵}。它们都可以用简单的策略节省大量存储空间。

\slidefig{assets/3.4-p02.jpg}{0.5}{课件：知识总览}

王道书第 100
页开头说：在数据结构中，关注的重点并非矩阵本身的数学性质及其运算，而是\uln{如何以最小的内存空间高效存储矩阵，并支持对元素的便捷访问}。

这一节所有的下标计算题，老师用的都是同一个办法：\hlred{数出这个元素前面存了多少个元素}，再看数组下标从
0 还是从 1
开始。老师反复说这些公式\uln{不要记}，考试一定会改条件，要能在考场上临时推出来。下面每种矩阵都画了一张“数个数”的图：橙色格子就是要数的那些。

\hypertarget{ux4e00ux4e00ux7ef4ux6570ux7ec4ux548cux4e8cux7ef4ux6570ux7ec4ux7684ux5b58ux50a8ux7ed3ux6784}{%
\section{一、一维数组和二维数组的存储结构}\label{ux4e00ux4e00ux7ef4ux6570ux7ec4ux548cux4e8cux7ef4ux6570ux7ec4ux7684ux5b58ux50a8ux7ed3ux6784}}

\tocback

\Needspace{4\baselineskip}

\cuenote{一维数组 a{[}i{]} 的地址怎么算？二维数组行优先、列优先的地址各怎么算？}\textbf{一维数组}。C
语言里 \texttt{ElemType\ a{[}10{]};}
定义一个一维数组。各数据元素类型相同，所以\uln{大小相同}，并且在内存中\uln{连续存放}。只要知道起始地址
LOC，就能算出任何一个下标对应的地址：

\[a[i] \text{ 的存放地址} = \text{LOC} + i \times \text{sizeof(ElemType)} \quad (0 \le i < 10)\]

\hlred{除非题目特别说明，数组下标默认从 0 开始}。如果从 1 开始，这里的
\(i\) 要变成 \(i-1\)，做题时一定要看清楚。

\slidefig{assets/3.4-p03.jpg}{0.5}{课件：一维数组的存储结构}

\textbf{二维数组}。\texttt{ElemType\ b{[}2{]}{[}4{]};} 定义一个 2 行 4
列的二维数组。逻辑视角上它有行有列，但内存空间是线性的，所以要把它“拉成”一条线，有两种存法：

\begin{itemize}
\tightlist
\item
  \textbf{行优先}：一行一行地存，先存第一行，紧接着存第二行\ldots\ldots{}
\item
  \textbf{列优先}：一列一列地存。
\end{itemize}

规定行优先或列优先，本质目的是把非线性的二维数组拉成线性的形状。这样有规律地存放之后，二维数组也能\uln{随机存取}：给出行号、列号，计算机就能直接算出地址。算法很简单，\uln{算出它前面有多少个元素}，乘以每个元素的大小，再加起始地址：

\begin{itemize}
\tightlist
\item
  M 行 N 列的 \texttt{b{[}M{]}{[}N{]}}，\textbf{行优先}：前面有 \(i\)
  整行，每行 \(N\) 个，本行前面还有 \(j\) 个，
  \[b[i][j] \text{ 的存放地址} = \text{LOC} + (i \times N + j) \times \text{sizeof(ElemType)}\]
\item
  \textbf{列优先}：前面有 \(j\) 整列，每列 \(M\) 个，本列前面还有 \(i\)
  个，
  \[b[i][j] \text{ 的存放地址} = \text{LOC} + (j \times M + i) \times \text{sizeof(ElemType)}\]
\end{itemize}

老师说这只是一个很简单的数学问题，这里主要理解什么叫行优先、什么叫列优先。

\slidefig{assets/3.4-p05.jpg}{0.5}{课件：二维数组的行优先存储}

\slidefig{assets/3.4-p06.jpg}{0.5}{课件：二维数组的列优先存储}

下面用课件的 \texttt{b{[}2{]}{[}4{]}} 找
\texttt{b{[}1{]}{[}2{]}}，把“数前面有几个”画出来（我画的）：

\slidefig{assets/3.4-2d.png}{0.85}{二维数组 b{[}2{]}{[}4{]} 中 b{[}1{]}{[}2{]} 的位置：行优先前面 6
个，列优先前面 5 个}

\begin{bookbox}

王道书第 100--101 页 §3.4.1 数组的定义：数组是由
\(n\)（\(n \ge 1\)）个相同类型的数据元素组成的有限序列。每个元素称为\textbf{数组元素}，其在线性序列中的位置由下标标识，下标的取值范围称为\textbf{维界}。

\textbf{数组与线性表的关系}：\uln{数组是线性表的推广}。一维数组可视为一个线性表；二维数组可视为其元素是定长一维数组的线性表；更高维数组以此类推。数组一旦定义，其维数与维界即固定不变。因此，除初始化和销毁外，\uln{数组仅支持存取元素和修改元素两种基本操作}。

§3.4.2
数组的存储结构：逻辑上的数组通常\uln{映射为内存中一段连续的存储空间}。一维数组
\(A[0\ldots n-1]\)，设每个元素占 \(L\) 个存储单元，则
\(\text{LOC}(a_i) = \text{LOC}(a_0) + iL \ (0 \le i < n)\)。

考点追踪（二维数组按行优先存储的下标对应关系，2021）：多维数组在内存中需按某种顺序展开为一维，常见方式有按行优先和按列优先。\textbf{按行优先}：先行后列，先存储行号较小的元素；同一行内，先存储列号较小的元素。设二维数组行下标与列下标范围分别为
\([0,h_1]\) 与 \([0,h_2]\)，则

\[\text{LOC}(a_{i,j}) = \text{LOC}(a_{0,0}) + [i(h_2+1) + j]L\]

\textbf{按列优先}：先列后行，先存储列号较小的元素；同一列内，先存储行号较小的元素。

\[\text{LOC}(a_{i,j}) = \text{LOC}(a_{0,0}) + [j(h_1+1) + i]L\]

\end{bookbox}

\begin{defbox}[：名词统一]

课件 \texttt{b{[}M{]}{[}N{]}}（M 行 N 列）和书上的维界
\([0,h_1]\)、\([0,h_2]\)：\(M = h_1+1\)，\(N = h_2+1\)，所以课件的
\(i \times N + j\) 就是书上的 \(i(h_2+1)+j\)。课件的 sizeof(ElemType) ＝
书上的 \(L\)；课件的起始地址 LOC ＝ 书上的 \(\text{LOC}(a_0)\) 或
\(\text{LOC}(a_{0,0})\)。

\end{defbox}

\hypertarget{ux4e8cux666eux901aux77e9ux9635ux7684ux5b58ux50a8}{%
\section{二、普通矩阵的存储}\label{ux4e8cux666eux901aux77e9ux9635ux7684ux5b58ux50a8}}

\tocback

\Needspace{4\baselineskip}

\cuenote{矩阵的行号列号和数组下标，默认各从几开始？}普通矩阵最容易想到用二维数组存储。要注意：\hlred{描述矩阵元素时，行号、列号通常从 1 开始；描述数组时，下标通常从 0 开始}。这是默认情况，做题时要仔细看题目给的条件，也有可能矩阵的行号、列号从
0 开始，这些都是比较容易错的地方。

普通矩阵可以用二维数组很方便地存储；而某些特殊矩阵，可以用一些巧妙的方法来压缩，节省存储空间。

\slidefig{assets/3.4-p07.jpg}{0.5}{课件：普通矩阵的存储}

\begin{bookbox}

王道书第 102 页“注意”：二维数组 \(A[n][n]\) 与
\(A[0\ldots n-1][0\ldots n-1]\) 写法等价，表示下标从 0 开始。若写作
\(A[1\ldots n][1\ldots n]\)，则表示下标从 1 开始。矩阵元素通常记为
\(a_{i,j}\)，其行号 \(i\) 和列号 \(j\) 通常从 1 开始。

王道书第 101 页
§3.4.3：\textbf{压缩存储}：为多个值相同的元素只分配一个存储空间，对零元素不分配存储空间。\textbf{特殊矩阵}：指具有大量相同元素或零元素，且这些元素分布具有规律性的矩阵。常见的特殊矩阵有对称矩阵、上（下）三角矩阵、对角矩阵等。\textbf{特殊矩阵的压缩存储方法}：通过分析矩阵中相同元素的分布规律，仅存储一份实际数据，其余的可通过下标映射访问，从而将原本冗余的数据压缩到一个共享的空间中，显著节省内存。

\end{bookbox}

\hypertarget{ux4e09ux5bf9ux79f0ux77e9ux9635}{%
\section{三、对称矩阵}\label{ux4e09ux5bf9ux79f0ux77e9ux9635}}

\tocback

\Needspace{4\baselineskip}

\cuenote{对称矩阵只存哪一部分？一维数组要多大？\(a_{i,j}\) 映射到哪个
\(k\)？\(i<j\) 怎么办？}\textbf{定义}。首先它是一个\textbf{方阵}（\(n\)
行 \(n\) 列），并且对任意元素都有 \(a_{i,j} = a_{j,i}\)，例如
\(a_{1,2}=a_{2,1}\)、\(a_{1,3}=a_{3,1}\)。元素的值关于主对角线对称（\(i=j\)
时式子显然也成立）。三个区域：

\begin{itemize}
\tightlist
\item
  \(i = j\)：\textbf{主对角线}；
\item
  \(i < j\)：\textbf{上三角区}；
\item
  \(i > j\)：\textbf{下三角区}。
\end{itemize}

\slidefig{assets/3.4-p08.jpg}{0.5}{课件：对称矩阵}

\textbf{压缩策略}。两个三角区的数据完全相同，所以\uln{只存主对角线＋下三角区}（也可以只存主对角线＋上三角区）。下面以前者为例，按\textbf{行优先}把这些元素依次存入一维数组
B。

\textbf{问题 1：数组要多大？} 第 1 行存 1 个，第 2 行 2
个\ldots\ldots 第 \(n\) 行 \(n\) 个，共

\[1 + 2 + \cdots + n = \frac{n(n+1)}{2}\]

个，所以 B 的最后一个元素下标是 \(\dfrac{n(n+1)}{2} - 1\)（减 1
是因为下标从 0 开始）。

\textbf{问题 2：怎么用矩阵下标访问？}
站在程序员的角度，存起来是为了用，用的时候肯定希望看到的是矩阵，而不是一维数组。所以要自己实现一个\textbf{映射函数}：矩阵下标
\(\to\)
一维数组下标。这是压缩存储之后要实现的很重要的东西，也是\hlred{考研最喜欢考的点}。

\begin{itemize}
\tightlist
\item
  \(i \ge j\)（下三角区和主对角线）：按行优先，\(a_{i,j}\) 前面有
  \(i-1\) 行，共 \(1+2+\cdots+(i-1)\) 个；列号 \(j\) 反映它是本行第
  \(j\) 个。所以它是第
  \([1+2+\cdots+(i-1)] + j = \dfrac{i(i-1)}{2} + j\) 个元素，数组下标从
  0 开始，还要减 1： \[k = \frac{i(i-1)}{2} + j - 1\]
\item
  \(i < j\)（上三角区）：这部分没有实际存进数组，但对称矩阵有
  \(a_{i,j} = a_{j,i}\)
  这个美好的特性，所以\uln{访问 \(a_{i,j}\) 就转换成访问 \(a_{j,i}\)}，又回到了行号大于列号的情况：
  \[k = \frac{j(j-1)}{2} + i - 1\]
\end{itemize}

老师强调：有的题目数组下标从 1 开始，那样就\hlred{不需要再减这个 1}。

\slidefig{assets/3.4-p12.jpg}{0.5}{课件：对称矩阵的映射函数（行优先）}

\textbf{出题方法}。实际考试肯定不会考原题，会改一些细微的条件：存上三角还是下三角？行优先还是列优先？矩阵元素下标从
0 还是 1 开始？数组下标从 0 还是 1
开始？老师举了\textbf{列优先}存主对角线＋下三角区的例子，思路一样，算
\(a_{i,j}\) 是第几个：

\begin{itemize}
\tightlist
\item
  第 \(j\) 列之前有第 1 到第 \(j-1\) 列：第 1 列 \(n\) 个，第 2 列
  \(n-1\) 个\ldots\ldots 第 \(j-1\) 列 \(n-j+2\) 个，等差数列求和得
  \(\dfrac{(j-1)(2n-j+2)}{2}\) 个；
\item
  同一列中它前面还有 \(i - j\) 个（行号减列号，例如 \(a_{n,3}\) 前面有
  \(n-3\) 个）；
\item
  所以它是第 \(\dfrac{(j-1)(2n-j+2)}{2} + (i-j) + 1\)
  个，再根据数组下标从 0 还是 1 开始转换成 \(k\)。
\end{itemize}

\slidefig{assets/3.4-p13.jpg}{0.5}{课件：对称矩阵按列优先存储，以及出题方法}

老师说除此之外还可以存主对角线＋上三角区，做题方法都很简单。课件只推了
①②，③④ 没推。下面把四种情况放在一张图里（我画的），\(n=4\)：

\slidefig{assets/3.4-sym4.png}{0.95}{对称矩阵四种存法：数格子里的下标。① 和 ④ 编号互为转置，② 和 ③ 也是}

\begin{addbox}

\textbf{补全课件没推的 ③④，并说明为什么只要会 ①② 就够了}。

\emph{③ 上三角＋主对角线，行优先}（\(i \le j\)）：第 \(i\) 行之前，第 1
行 \(n\) 个，第 2 行 \(n-1\) 个\ldots\ldots 第 \(i-1\) 行 \(n-i+2\)
个，共 \(\dfrac{(i-1)(2n-i+2)}{2}\) 个；本行从 \(a_{i,i}\)
开始存，它前面还有 \(j-i\) 个。所以
\(k = \dfrac{(i-1)(2n-i+2)}{2} + j - i\)。\(i > j\) 时访问 \(a_{j,i}\)。

\emph{④ 上三角＋主对角线，列优先}（\(i \le j\)）：第 \(j\) 列之前有
\(1+2+\cdots+(j-1)\) 个，本列从 \(a_{1,j}\) 开始存，它是第 \(i\)
个。所以 \(k = \dfrac{j(j-1)}{2} + i - 1\)。

\emph{为什么 ③ 和 ② 是同一个公式（\(i\)、\(j\) 互换）}：③ 第 \(r\)
行存的是 \(a_{r,r}, a_{r,r+1}, \ldots, a_{r,n}\)，② 第 \(r\) 列存的是
\(a_{r,r}, a_{r+1,r}, \ldots, a_{n,r}\)。把 ② 的每个元素 \(a_{p,q}\)
换成 \(a_{q,p}\)，正好得到 ③ 的同一位置的元素，次序也一样。所以 ③ 中
\(a_{i,j}\) 的位置 ＝ ② 中 \(a_{j,i}\) 的位置，把 ② 的公式里
\(i\)、\(j\) 互换就是 ③ 的公式。同理，④ 中 \(a_{i,j}\) 的位置 ＝ ① 中
\(a_{j,i}\) 的位置。图里能直接看出来：① 中 \(a_{3,2}\) 是 4，④ 中
\(a_{2,3}\) 也是 4；② 中 \(a_{3,2}\) 是 5，③ 中 \(a_{2,3}\) 也是 5。

所以四种情况只有两种本质不同的数法：\uln{每段从 1 个开始变多}（①④，前面共
\(\frac{m(m-1)}{2}\) 个），\uln{每段从 \(n\) 个开始变少}（②③，前面共
\(\frac{m(2n-m+1)}{2}\) 个，\(m\) 是前面的段数）。

\emph{数组下标从 1 开始}：上面每个 \(k\) 都加 1。\emph{矩阵下标从 0
开始}：先把 \(i\)、\(j\) 各加 1，再代入公式。

\end{addbox}

\begin{bookbox}

王道书第 101--102
页：考点追踪（对称矩阵压缩存储的下标对应关系，2018、2020）。若 \(n\)
阶矩阵 \(A\) 中任意元素 \(a_{i,j}\) 均满足
\(a_{i,j} = a_{j,i}\)（\(1 \le i, j \le n\)），则称其为对称矩阵。其元素可分为三部分：上三角区、主对角线和下三角区。由于上三角区与下三角区完全对称，若仍用二维数组存储，将近一半空间被浪费。为此可将
\(n\) 阶对称矩阵 \(A\) 压缩存储于一维数组 \(B[n(n+1)/2]\)
中，通常仅存放下三角部分（含主对角线）。

对于元素 \(a_{i,j}\)（\(i \ge j\)），其在数组 B
中的位置由其前方的元素个数决定：第 1 行 1 个，第 2 行 2
个\ldots\ldots 第 \(i-1\) 行 \(i-1\) 个，第 \(i\) 行 \(j-1\)
个（\(a_{i,1}, \ldots, a_{i,j-1}\)）。因此
\(k = 1+2+\cdots+(i-1)+j-1 = i(i-1)/2 + j - 1\)（数组下标从 0 开始），

\[k = \begin{cases} \dfrac{i(i-1)}{2} + j - 1, & i \ge j \text{（下三角区和主对角线元素）} \\[2mm] \dfrac{j(j-1)}{2} + i - 1, & i < j \text{（上三角区元素 } a_{i,j} = a_{j,i}\text{）} \end{cases}\]

若数组下标从 1 开始，则可采用类似的推导方法，请读者自行思考。

\end{bookbox}

\hypertarget{ux56dbux4e09ux89d2ux77e9ux9635}{%
\section{四、三角矩阵}\label{ux56dbux4e09ux89d2ux77e9ux9635}}

\tocback

\Needspace{4\baselineskip}

\cuenote{三角矩阵比同阶对称矩阵多存什么？常量区映射到哪里？上三角矩阵行优先的
\(k\) 怎么推？}\textbf{定义}。\textbf{下三角矩阵}：除了主对角线和下三角区，其余元素（上三角区）都相同；\textbf{上三角矩阵}：除了主对角线和上三角区，其余元素（下三角区）都相同。

\slidefig{assets/3.4-p14.jpg}{0.5}{课件：下三角矩阵和上三角矩阵}

\textbf{压缩策略}。常量区那一整块的值都相同，不需要重复存放那么多份，只需要重点存储不是常量的部分，方法和对称矩阵一样（行优先或列优先存入一维数组）；\uln{在最后再加一个位置存放这个常量 \(c\)}。所以一维数组长度是
\(\dfrac{n(n+1)}{2} + 1\)，比存储同阶对称矩阵\hlred{多一个存储单元}。

\textbf{下三角矩阵（行优先）}：\(i \ge j\)
时和对称矩阵的式子一模一样；唯一不同的是
\(i < j\)（上三角区）的元素都是常量
\(c\)，\uln{全部映射到一维数组的最后一个位置}：

\[k = \begin{cases} \dfrac{i(i-1)}{2} + j - 1, & i \ge j \\[2mm] \dfrac{n(n+1)}{2}, & i < j \end{cases}\]

\slidefig{assets/3.4-p15.jpg}{0.5}{课件：下三角矩阵的压缩存储}

\textbf{上三角矩阵（行优先）}：同样考虑 \(a_{i,j}\) 前面有几个元素。第
\(i\) 行之前，第 1 行 \(n\) 个，第 2 行 \(n-1\) 个\ldots\ldots 第
\(i-1\) 行 \(n-i+2\) 个，等差数列求和；第 \(i\) 行中 \(a_{i,j}\)
前面还有 \(j - i\) 个。所以

\[k = \begin{cases} \dfrac{(i-1)(2n-i+2)}{2} + (j-i), & i \le j \\[2mm] \dfrac{n(n+1)}{2}, & i > j \text{（下三角区，映射到常量的位置）} \end{cases}\]

\slidefig{assets/3.4-p16.jpg}{0.5}{课件：上三角矩阵的压缩存储}

两种三角矩阵的数法画在一起（我画的），\(n=4\)，常量 \(c\) 存在
\(B[10]\)：

\slidefig{assets/3.4-tri.png}{0.9}{三角矩阵：数前面有几个，常量区一律映射到最后一个位置}

课件第 20 页在下三角矩阵旁边留了一个思考题：“如何由 \(k\) 推出
\(i\)、\(j\)？”，正课视频里老师没有讲，下面补上。

\begin{addbox}

\textbf{补：下三角矩阵（行优先），已知 \(k\) 求
\(i\)、\(j\)}（思路和第五节三对角矩阵老师讲的“刚好”一样）。

\(k = \frac{n(n+1)}{2}\) 时是常量 \(c\)。否则 \(B[k]\) 前面有 \(k\)
个元素。前 \(i-1\) 行共 \(\frac{i(i-1)}{2}\) 个，前 \(i\) 行共
\(\frac{i(i+1)}{2}\) 个。\(B[k]\) 在第 \(i\) 行，当且仅当

\[\frac{i(i-1)}{2} \le k < \frac{i(i+1)}{2}.\]

左边：\(i^2 - i - 2k \le 0\)，即
\(i \le \frac{1+\sqrt{8k+1}}{2}\)；右边：\(i^2 + i - 2k > 0\)，即
\(i > \frac{-1+\sqrt{8k+1}}{2} = \frac{1+\sqrt{8k+1}}{2} - 1\)。记
\(x = \frac{1+\sqrt{8k+1}}{2}\)，\(i\) 是区间 \((x-1, x]\)
里的整数，这个区间长度为 1，里面恰好有一个整数，就是
\(\lfloor x \rfloor\)：

\[i = \left\lfloor \frac{1+\sqrt{8k+1}}{2} \right\rfloor, \qquad j = k - \frac{i(i-1)}{2} + 1.\]

\emph{例子}（对照上图）：\(k=8\)，\(\sqrt{65} \approx 8.06\)，\(x \approx 4.53\)，\(i = 4\)；\(j = 8 - 6 + 1 = 3\)，是
\(a_{4,3}\)。\(k=6\)，\(\sqrt{49} = 7\)，\(x = 4\)，\(i = 4\)，\(j = 1\)，是
\(a_{4,1}\)，正好是第 4 行的第一个。

\end{addbox}

\begin{bookbox}

王道书第 102--103
页：下三角矩阵中，上三角区的所有元素均为同一常量。其存储思想与对称矩阵的类似，但是需要额外存储该常量一次。因此，可以将
\(n\) 阶下三角矩阵 \(A\) 压缩存储在 \(B[n(n+1)/2+1]\) 中，\(k\)
的式子和上面相同（图 3.21：\(B[0]\) 起依次是第 1 行、第 2
行\ldots\ldots 第 \(n\) 行，\(B[n(n+1)/2]\) 是常数项）。

考点追踪（上三角矩阵采用行优先存储的应用，2011）：上三角矩阵中，下三角区所有元素均为同一常量。只需存储主对角线、上三角区上元素及该常量一次，同样将其压缩存储在
\(B[n(n+1)/2+1]\) 中。对于元素 \(a_{i,j}\)（\(i \le j\)），其在数组 B
中前方的元素个数为：第 1 行 \(n\) 个，第 2 行 \(n-1\) 个\ldots\ldots 第
\(i-1\) 行 \(n-i+2\) 个，第 \(i\) 行 \(j-i\) 个。故
\(k = n+(n-1)+\cdots+(n-i+2)+(j-i+1)-1 = (i-1)(2n-i+2)/2+(j-i)\)。

以上推导均假设数组下标从 0 开始。若下标指定从 1
开始，则要相应地调整映射关系。

\end{bookbox}

\hypertarget{ux4e94ux4e09ux5bf9ux89d2ux77e9ux9635}{%
\section{五、三对角矩阵}\label{ux4e94ux4e09ux5bf9ux89d2ux77e9ux9635}}

\tocback

\Needspace{4\baselineskip}

\cuenote{三对角矩阵哪些元素可以非零？要存几个？\(k = 2i+j-3\) 怎么来？已知 \(k\)
怎么求 \(i\)、\(j\)？}\textbf{定义}。\textbf{三对角矩阵}又称\textbf{带状矩阵}：当
\(|i - j| > 1\)
时，\(a_{i,j} = 0\)（\(1 \le i, j \le n\)）。形象地理解：主对角线上
\(|i-j| = 0\)；往右（或左、上、下）走一格，差值变为
1，可以不为零；再走一格差值变为
2，肯定是零。所以\uln{主对角线上的元素，以及和它们上下左右相邻的元素可以非零}，再往外都是零。

\textbf{压缩策略}。只存带状部分，按行优先（或列优先）依次存入一维数组。除了第一行和最后一行只有
2 个元素，其他每行都有 3 个，所以共

\[3n - 2\]

个元素，最后一个的数组下标是 \(3n - 3\)。

\textbf{正向：\(a_{i,j} \to k\)}（只讨论 \(|i - j| \le 1\)，差值大于 1
的元素肯定是 0，不用来数组里找）。每行 3 个，第一行少一个，所以前
\(i-1\) 行共 \(3(i-1) - 1\) 个；\(a_{i,j}\) 是第 \(i\) 行的第
\(j - i + 2\) 个。加起来，\(a_{i,j}\) 是第 \(2i + j - 2\) 个元素，下标从
0 开始：

\[k = 2i + j - 3\]

\slidefig{assets/3.4-p17.jpg}{0.5}{课件：三对角矩阵的压缩存储}

\textbf{反向：\(k \to i, j\)}。\(B[k]\) 是第 \(k+1\) 个元素。前 \(i-1\)
行共 \(3(i-1)-1\) 个，前 \(i\) 行共 \(3i - 1\) 个，它在第 \(i\)
行就要满足

\[3(i-1) - 1 < k + 1 \le 3i - 1.\]

从右边解出 \(i \ge \dfrac{k+2}{3}\)。\(i\) 是行号，一定是整数，要让
\(i\) \uln{刚好}大于等于右边这个值，\hlred{向上取整}即可：

\[i = \left\lceil \frac{k+2}{3} \right\rceil, \qquad j = k - 2i + 3 \ \text{（由 } k = 2i+j-3\text{）}.\]

\textbf{王道书的算法}：把 \(k\) 理解为它前面共有 \(k\) 个元素，那么
\(3(i-1) - 1 \le k < 3i - 1\)。从左边解出
\(i \le \dfrac{k+1}{3} + 1\)，要 \(i\)
刚好小于等于这个值，\hlred{向下取整}：\(i = \left\lfloor \dfrac{k+1}{3} + 1 \right\rfloor\)。老师说书上的式子看起来不一样，计算结果完全一样，“大家可以自己带进去验证”，“背后的逻辑留给大家自己去琢磨”，要注意体会不等式里“刚好大于等于”“刚好小于等于”怎么处理。

\slidefig{assets/3.4-p18.jpg}{0.5}{课件：已知 k 求 i（课件的算法和王道书的算法）}

\slidefig{assets/3.4-p19.jpg}{0.5}{课件：已知 k 求 i、j 的结果}

下面把 \(n=5\)
的三对角矩阵按下标编好号（我画的）。下面一条是一维数组，每行占一段：

\slidefig{assets/3.4-tridiag.png}{0.9}{三对角矩阵：\(a_{3,4}\) ↔ k＝7。第 3 行占 k＝5～7，所以 k＝7 在第 3 行}

\begin{addbox}

\textbf{证明老师留的两件事}。

\emph{1. \(i = \lceil \frac{k+2}{3} \rceil\) 和
\(i = \lfloor \frac{k+1}{3} \rfloor + 1\) 是同一个数。}
两种算法都是在找满足一对不等式的整数 \(i\)，只要证明对整数 \(i\)
这两对不等式等价。对整数 \(i\)、\(k\)：

\begin{itemize}
\tightlist
\item
  \(i \ge \frac{k+2}{3} \iff 3i \ge k+2 \iff 3i > k+1 \iff i > \frac{k+1}{3}\)；
\item
  老师的另一半 \(3(i-1)-1 < k+1\)，即
  \(3i < k+5 \iff 3i \le k+4 \iff i \le \frac{k+1}{3}+1\)，正是书上的条件。
\end{itemize}

所以两种算法找的是同一个整数：它在区间
\(\left(\frac{k+1}{3}, \frac{k+1}{3}+1\right]\) 里，区间长度为
1，恰好有一个整数，就是 \(\lfloor \frac{k+1}{3} \rfloor + 1\)，也等于
\(\lceil \frac{k+2}{3} \rceil\)。（“\(>\)”能换成“\(\ge\) 再加
1”，靠的就是 \(i\)、\(k\) 都是整数，这也是老师说的“刚好”。）

\emph{2. 书上为什么“\(\le\)”“\(<\)”的位置和老师的相反。}
老师数的是“它是第 \(k+1\) 个”，书上数的是“它前面有 \(k\) 个”，差 1。前
\(i-1\) 行的个数 \(3(i-1)-1\)：老师要求严格小于 \(k+1\)，书上要求
\(\le k\)，对整数是同一个条件；前 \(i\) 行的个数 \(3i-1\)：老师要求
\(\ge k+1\)，书上要求 \(> k\)，也是同一个条件。

\emph{3. 第一行的公式也对，但有两处“差 1”正好抵消。} 第 1 行只有
\(a_{1,1}\)、\(a_{1,2}\)，\(a_{1,j}\) 实际是本行第 \(j\) 个（不是
\(j - i + 2 = j+1\) 个），而“前 0 行共 \(3\times0-1 = -1\) 个”也少算了
1。两处相加，\(2i+j-2 = j\) 仍是正确的序号，所以 \(k = 2i+j-3\) 对第 1
行也成立：\(a_{1,1}\) 得 0，\(a_{1,2}\) 得 1。

\emph{例子}（对照上图）：\(k = 7\)，老师的算法
\(\lceil 9/3 \rceil = 3\)，书上的算法
\(\lfloor 8/3 \rfloor + 1 = 2 + 1 = 3\)，同一行；\(j = 7 - 6 + 3 = 4\)，是
\(a_{3,4}\)。图里第 3 行占 \(k = 5\sim7\)，前 2 行共 \(3\times2-1 = 5\)
个，正好是 \(k=0\sim4\)。

\end{addbox}

\begin{bookbox}

王道书第 103 页：若 \(n\) 阶矩阵 \(A\) 中的任意元素 \(a_{i,j}\)，当
\(|i-j| > 1\) 时均有
\(a_{i,j} = 0\)（\(1 \le i, j \le n\)），则称为三对角矩阵。非零元素仅集中在以主对角线为中心的
3 条对角线的区域。可将三对角上的元素按行优先顺序存入一维数组 B，且
\(a_{1,1}\) 存于 \(B[0]\)。

考点追踪（三对角矩阵压缩存储的下标对应关系，2016）：三对角上的元素
\(a_{i,j}\)（\(1 \le i, j \le n\)，\(|i-j| \le 1\)）在一维数组 B
中的下标 \(k = 2i+j-3\)。反之，已知元素存于 \(B[k]\)
中时，\(i = \lfloor (k+1)/3 + 1 \rfloor\)，\(j = k - 2i + 3\)。例如，当
\(k = 0\)
时，\(i = \lfloor (0+1)/3 + 1 \rfloor = 1\)，\(j = 0 - 2 + 3 = 1\)，存放的是
\(a_{1,1}\)；当 \(k = 2\) 时，\(i = 2\)，\(j = 1\)，存放的是
\(a_{2,1}\)；当 \(k = 4\) 时，\(i = 2\)，\(j = 3\)，存放的是
\(a_{2,3}\)。

\end{bookbox}

\hypertarget{ux516dux7a00ux758fux77e9ux9635}{%
\section{六、稀疏矩阵}\label{ux516dux7a00ux758fux77e9ux9635}}

\tocback

\Needspace{4\baselineskip}

\cuenote{稀疏矩阵的两种压缩存储方法是什么？三元组存储后失去了什么？十字链表的结点有哪几个域？}\textbf{定义}。\textbf{稀疏矩阵}：\uln{非零元素的个数远远少于矩阵元素的个数}。什么叫“远远少于”并没有固定的界限，老师说作为程序员大家应该对“稀疏”这个词都有很深刻的理解，“摸摸自己的头顶就知道了”。如果定义一个和它尺寸相当的二维数组，很多空间都存的是
0，没有必要。

\textbf{方法一：顺序存储三元组
⟨行，列，值⟩}。每个非零元素用一个三元组表示，例如课件中 4 在第 1 行第 3
列，三元组是 ⟨1, 3, 4⟩；5 在第 1 行第 6 列，是 ⟨1, 6,
5⟩。顺序存储的做法：定义一个 struct，里面有
\texttt{i}、\texttt{j}、\texttt{value} 三个字段（一个 struct
对应表中的一行），再定义一个这种 struct
的一维数组。用这种方式存储，要访问某个元素，只能\uln{顺序依次扫描这些三元组}，也就是\hlred{失去了随机存取的特性}。

\slidefig{assets/3.4-p21.jpg}{0.5}{课件：稀疏矩阵的顺序存储（三元组）}

\textbf{方法二：十字链表法}（老师说课本当中没有详细讲）。“一图胜千言”：

\begin{itemize}
\tightlist
\item
  一个指针数组
  \textbf{向下域}（down），每个指针对应稀疏矩阵的一列，\uln{指向第 \(j\) 列的第一个元素}；
\item
  一个指针数组
  \textbf{向右域}（right），每个指针对应一行，\uln{指向第 \(i\) 行的第一个元素}；
\item
  每个非零元素对应一个结点：行、列、值，以及两个指针，\uln{一个指向同列的下一个元素，一个指向同行的下一个元素}。
\end{itemize}

例如第 1 行的第一个非零元素是 4，第 1 个向右域指针指向结点 ⟨1, 3,
4⟩，这个结点的指针再指向同一行的下一个非零元素 ⟨1, 6, 5⟩。第 2
列从上往下第一个非零元素是 3，第 2 个向下域指针指向结点 ⟨2, 2,
3⟩，它再指向同一列的下一个非零元素 ⟨4, 2, 2⟩。

\slidefig{assets/3.4-p22.jpg}{0.6}{课件：稀疏矩阵的链式存储（十字链表法）}

\begin{bookbox}

王道书第 103--104 页 §3.4.4：若矩阵中非零元素的个数 \(t\)
相对于总元素个数 \(s\) 来说非常少，即
\(t \ll s\)，则称该矩阵为稀疏矩阵。例如，一个 \(100\times100\)
的矩阵中只有不到 100 个非零元素。

考点追踪（存储稀疏矩阵需要保存的信息，2023）：为避免空间浪费，稀疏矩阵通常仅存储非零元素。然而，由于非零元素的分布通常是无规律的，仅存储其值是不够的，还需记录它们所在的行和列。为此，将每个非零元素及其对应的行列位置组合成一个\textbf{三元组}（行标
\(i\)，列标 \(j\)，值
\(a_{i,j}\)）。这些三元组可以按某种顺序排列成线性表进行存储。\uln{稀疏矩阵压缩存储后便失去了随机存取特性}。

考点追踪（稀疏矩阵压缩存储的方式，2017）：稀疏矩阵的三元组表可通过多种方式存储，常见的有：\textbf{数组存储}将所有三元组按某种顺序存储在一维数组中，这种方式简单直接，但插入和删除操作的效率较低；\textbf{十字链表存储}（见
6.2
节）通过链表的方式组织三元组，适用于频繁插入和删除操作的场景。无论采用哪种方式，都需要\hlred{额外保存稀疏矩阵的行数、列数和非零元素的个数}，以便支持后续的各种操作。

\end{bookbox}

\slidefig{assets/3.4-book-fig3.26.png}{0.5}{王道书图 3.26：稀疏矩阵及其对应的三元组（注意这里行列从 0 开始）}

\begin{defbox}[：名词统一]

课件“三对角矩阵，又称带状矩阵”＝书上“三对角矩阵”。课件“顺序存储三元组”＝书上“数组存储”；课件“链式存储：十字链表法”＝书上“十字链表存储”。课件稀疏矩阵的三元组\uln{行、列从 1 开始}（课件注明），书上图
3.26 的三元组行、列\uln{从 0 开始}，做题看清题目。课件“向下域 down /
向右域 right”，书上这里没有讲，到第 6 章 6.2 节图的十字链表再学。

\end{defbox}

\hypertarget{ux4e03ux77e5ux8bc6ux56deux987e}{%
\section{七、知识回顾}\label{ux4e03ux77e5ux8bc6ux56deux987e}}

\tocback

\slidefig{assets/3.4-p23.jpg}{0.5}{课件：知识回顾与重要考点}

\slidefig{assets/3.4-p24.jpg}{0.5}{课件：常见考题和易错点}

老师说这部分内容有点多，但难度不大，考题相对固定。\hlred{题目中可能会有坑}：存上三角还是下三角？行优先还是列优先？矩阵元素下标从
0 还是 1 开始？数组下标从 0 还是 1
开始？课件最后还问了一句：“是否忘了，等差数列求和？”

\begin{summarybox}

\textbf{数组}：元素大小相同、连续存放，可随机存取。一维
\(a[i]\)：\(\text{LOC} + i \times \text{sizeof}\)；二维 \(b[M][N]\)
行优先 \(\text{LOC} + (iN + j) \times \text{sizeof}\)，列优先
\(\text{LOC} + (jM + i) \times \text{sizeof}\)。数组下标默认从 0
开始，矩阵行列号默认从 1 开始。

\textbf{统一做法}：\uln{数出前面有几个元素}（等差数列求和），再看数组下标从
0 还是 1 开始。

\textbf{对称矩阵}：\(a_{i,j} = a_{j,i}\)；只存主对角线＋下（上）三角区，\(\frac{n(n+1)}{2}\)
个。下三角行优先：\(i \ge j\) 时
\(k = \frac{i(i-1)}{2} + j - 1\)，\(i<j\) 时访问 \(a_{j,i}\)。

\textbf{三角矩阵}：常量区只存一次，放在最后，共 \(\frac{n(n+1)}{2} + 1\)
个；常量区元素映射到
\(k = \frac{n(n+1)}{2}\)。上三角行优先：\(k = \frac{(i-1)(2n-i+2)}{2} + (j-i)\)。

\textbf{三对角矩阵（带状矩阵）}：\(|i-j|>1\) 时为 0，共 \(3n-2\)
个；\(k = 2i + j - 3\)；反向 \(i = \lceil \frac{k+2}{3} \rceil\)（＝书上
\(\lfloor \frac{k+1}{3} \rfloor + 1\)），\(j = k - 2i + 3\)。

\textbf{稀疏矩阵}：非零元素远少于总数。① 顺序存储三元组
⟨行，列，值⟩，\hlred{失去随机存取特性}；②
十字链表法。都要额外保存行数、列数和非零元素个数。

\textbf{王道书考点}：二维数组按行优先存储的下标（2021）；对称矩阵下标对应关系（2018、2020）；上三角矩阵行优先（2011）；三对角矩阵下标（2016）；存储稀疏矩阵要保存的信息（2023）；稀疏矩阵压缩存储的方式（2017）。

\textbf{自测}：① \(n\)
阶对称矩阵按列优先存下三角，\(a_{i,j}\)（\(i \ge j\)）的 \(k\)？②
上三角矩阵行优先，\(a_{2,4}\)（\(n=4\)）存在哪？常量 \(c\) 呢？③
三对角矩阵 \(k=10\)（\(n=5\)）是哪个元素？④
稀疏矩阵三元组存储后失去了什么特性？

\end{summarybox}
\end{document}
